\scanpage{7}
\mspage{5}

schéma fini étale sur $S$, alors on a des isomorphismes
\[
\operatorname{Hom}_S(\bar s,S')\xleftarrow{\sim}\operatorname{Hom}_S(\deleted{\unsure{\bar s}}\widetilde S,S')\xrightarrow{\sim}\operatorname{Hom}_S(\deleted{\unsure{\eta}}\bar\eta,S')
\]
d'où un isom canonique
\[
(2.3.1)\qquad\operatorname{Hom}_S(\bar s,S')\simeq\operatorname{Hom}_S(\bar\eta,S'),
\]
fonctoriel en $S'$ (isom. des \inserted{2} foncteurs ``fibre'' sur la cat. des rev. étales de $S$, définis resp. par les pts géom. $\bar\eta$ et $\bar s$). Faisant $S'=\mu_n$, $n$ \inserted{entier} inversible dans $k$, on trouve en particulier un isom
\[
(2.\unsure{1|3}.2)\qquad \mu_n(\bar s)={}_n\bar k^*\simeq\mu_n(\bar K)={}_n\bar K^*,
\]
d'où par passage à la limite sur les puissances du nb premier $\ell$, $\ell\ne\operatorname{car}.k$ :
\[
(2.3.3)\qquad T_\ell(\bar k^*)\simeq T_\ell(\bar K^*).
\]
\deleted{On notera.} Les deux membres, \deleted{qui} sont des modules libres de rang 1 sur $\mathbb Z_\ell$, \deleted{\unsure{que l'on désigne}} \inserted{qu'on identifie} désormais et désigne par $\mathbb Z_\ell(1)$. \deleted{Sa puissance} \inserted{Sa puiss. tensorielle} nième $(n\in\mathbb Z)$ est notée $\mathbb Z_\ell(n)$. Donc on a
\margin{$(2.3.\unsure{3|4})\qquad\mathbb Z_\ell(1)\simeq\cdots$}
\[
(2.3.4)\qquad \mathbb Z_\ell(n)\simeq\varprojlim_\nu\mu_{\ell^\nu}(\bar k^*)^{\otimes n}\simeq\varprojlim_\nu\mu_{\ell^\nu}(\bar K^*)^{\otimes n}.
\]
Les isomorphismes \deleted{\unsure{sont compatibles aux}} opérations de $\pi$ sur $\mu_{\unsure{n}}(\bar K^*)$ \inserted{et} \unsure{ou} celles de $\pi_0$ sur $\mu_n(\bar k^*)$ etc\ldots compte tenu de l'épimorphisme canonique $\pi\unsure{-|\longrightarrow}\pi_0$.
\margin{6}
